\documentclass[../DoAn.tex]{subfiles}
\usepackage{float}
\usepackage{tabularx}
\usepackage{array}
\usepackage{ragged2e}

\begin{document}

\section{Công nghệ và môi trường phát triển}
\label{section:4.1}

Trong quá trình phát triển hệ thống quản lý nhà hàng, nhóm đã lựa chọn các công nghệ phù hợp nhằm đảm bảo hiệu năng, khả năng mở rộng và dễ bảo trì.

\subsection{Công nghệ sử dụng}

\textbf{Frontend:}
\begin{itemize}
    \item \textbf{Vue.js 3:} Framework JavaScript theo mô hình component, hỗ trợ tái sử dụng giao diện và tối ưu hiệu năng.
    \item \textbf{Vite:} Công cụ build và dev server cho ứng dụng Vue, hỗ trợ tải module nhanh trong quá trình phát triển.
    \item \textbf{Element Plus:} Thư viện UI chính cho các màn hình khách hàng và một phần khu vực quản trị.
    \item \textbf{Ant Design Vue:} Bổ sung component UI cho các màn hình quản trị (bảng dữ liệu, form, thống kê).
    \item \textbf{Vue Router:} Quản lý điều hướng giữa các trang theo mô hình SPA (Single Page Application).
    \item \textbf{Axios:} Gửi yêu cầu HTTP tới REST API của backend.
    \item \textbf{ECharts:} Hiển thị biểu đồ trong module báo cáo doanh thu.
\end{itemize}

\textbf{Backend:}
\begin{itemize}
    \item \textbf{Node.js:} Môi trường chạy JavaScript phía server, xử lý bất đồng bộ hiệu quả.
    \item \textbf{Express.js:} Framework xây dựng RESTful API.
    \item \textbf{Sequelize:} ORM kết nối và thao tác với PostgreSQL.
    \item \textbf{JSON Web Token (JWT):} Xác thực phiên đăng nhập và phân quyền theo vai trò.
    \item \textbf{bcrypt:} Mã hóa mật khẩu trước khi lưu cơ sở dữ liệu.
    \item \textbf{WebSocket:} Đồng bộ trạng thái món ăn giữa bếp và nhân viên phục vụ theo từng chi nhánh.
\end{itemize}

\textbf{Cơ sở dữ liệu:}
\begin{itemize}
    \item \textbf{PostgreSQL:} Hệ quản trị cơ sở dữ liệu quan hệ, đảm bảo tính toàn vẹn dữ liệu và hỗ trợ truy vấn phức tạp phục vụ báo cáo.
\end{itemize}

\subsection{Môi trường phát triển}

\begin{itemize}
    \item \textbf{Hệ điều hành:} Windows 10/11
    \item \textbf{IDE:} Visual Studio Code
    \item \textbf{Trình duyệt:} Google Chrome
    \item \textbf{Công cụ quản lý phiên bản:} Git, GitHub
\end{itemize}

\subsection{Kiến trúc hệ thống}

Hệ thống được xây dựng theo mô hình \textbf{Client-Server}, trong đó:
\begin{itemize}
    \item \textbf{Frontend} chịu trách nhiệm hiển thị giao diện và gửi yêu cầu đến server.
    \item \textbf{Backend} xử lý logic nghiệp vụ và giao tiếp với cơ sở dữ liệu.
    \item \textbf{Database} lưu trữ dữ liệu người dùng, đặt bàn, đơn hàng, bàn, món ăn, thanh toán và đánh giá.
\end{itemize}

Ngoài ra, backend được tổ chức theo kiến trúc phân tầng:
\begin{itemize}
    \item \textbf{Routes:} Khai báo endpoint API và middleware (xác thực, phân quyền, ghi log).
    \item \textbf{Controller:} Nhận request từ client và trả response.
    \item \textbf{Service:} Xử lý logic nghiệp vụ.
    \item \textbf{Model:} Ánh xạ bảng dữ liệu và thao tác với PostgreSQL thông qua Sequelize.
\end{itemize}

Kênh \textbf{WebSocket} bổ sung cho REST API nhằm đẩy thông báo thay đổi trạng thái món ăn tới giao diện phục vụ và bếp trong cùng chi nhánh.

\subsection{Đánh giá lựa chọn công nghệ}

Các công nghệ được lựa chọn đều phổ biến, dễ triển khai và có cộng đồng hỗ trợ lớn.
Việc sử dụng Vue.js kết hợp Node.js giúp hệ thống đồng nhất về ngôn ngữ (JavaScript), giảm độ phức tạp khi phát triển.
PostgreSQL phù hợp với hệ thống quản lý quy mô vừa; JWT và bcrypt đáp ứng yêu cầu bảo mật cơ bản trong phạm vi đồ án.

\section{Cấu trúc dự án}
\label{section:4.2}

Dự án được tổ chức tách biệt thư mục \texttt{vue-fe/} (frontend) và \texttt{be/} (backend) nhằm dễ bảo trì, mở rộng và triển khai.

\subsection{Cấu trúc frontend}

Frontend được xây dựng bằng Vue.js~3 theo hướng component-based:

\begin{itemize}
    \item \textbf{src/views:} Các trang chính (đăng nhập, đặt bàn, menu, quản lý bàn, bếp, báo cáo\ldots).
    \item \textbf{src/components:} Component dùng chung (navbar, sidebar, form, bảng dữ liệu\ldots).
    \item \textbf{src/router:} Cấu hình điều hướng Vue Router.
    \item \textbf{src/layouts:} Khung giao diện khu vực quản trị.
    \item \textbf{src/utils:} Tiện ích xác thực, phạm vi chi nhánh, kết nối realtime.
    \item \textbf{src/config:} Cấu hình menu sidebar và nội dung trang tĩnh.
    \item \textbf{src/constants:} Hằng số trạng thái bàn, đơn hàng.
    \item \textbf{src/assets:} CSS, hình ảnh và tài nguyên giao diện.
\end{itemize}

Trạng thái phiên đăng nhập được lưu tại trình duyệt (token JWT); dự án \textbf{chưa} sử dụng thư viện quản lý state tập trung (Pinia/Vuex).

\subsection{Cấu trúc backend}

Backend được xây dựng bằng Node.js và Express theo kiến trúc phân tầng:

\begin{itemize}
    \item \textbf{routes:} Khai báo các API endpoint theo nhóm (người dùng, admin, public).
    \item \textbf{controllers:} Nhận request và trả response.
    \item \textbf{services:} Xử lý logic nghiệp vụ.
    \item \textbf{models:} Định nghĩa bảng dữ liệu và quan hệ (Sequelize).
    \item \textbf{middlewares:} Xác thực JWT, phân quyền, kiểm tra đầu vào, giới hạn tần suất, ghi nhật ký thao tác.
    \item \textbf{utils:} JWT, xuất hóa đơn PDF, tiện ích dùng chung.
    \item \textbf{seeders:} Dữ liệu mẫu phục vụ demo và kiểm thử.
    \item \textbf{realtimeHub.js:} Phát sự kiện WebSocket theo chi nhánh.
\end{itemize}

\subsection{Nguyên tắc tổ chức mã nguồn}

Mã nguồn được xây dựng theo hướng module hóa nhằm:
\begin{itemize}
    \item Dễ mở rộng chức năng trong tương lai.
    \item Giảm sự phụ thuộc giữa các module.
    \item Hỗ trợ tái sử dụng code.
    \item Tăng khả năng bảo trì và kiểm thử hệ thống.
\end{itemize}

\section{Thiết kế giao diện người dùng}
\label{section:4.3}

Giao diện người dùng được thiết kế theo định hướng trực quan, dễ sử dụng và phù hợp với từng nhóm tác nhân.
Do hệ thống phục vụ nhiều vai trò (khách hàng, nhân viên phục vụ, bếp, quản lý, admin),
giao diện được phân tách theo luồng nghiệp vụ và quyền truy cập, giúp thao tác nhanh và giảm nhầm lẫn
trong môi trường nhà hàng.

\subsection{Nguyên tắc thiết kế}

Các giao diện được xây dựng dựa trên các nguyên tắc sau:

\begin{itemize}
    \item \textbf{Đơn giản và dễ thao tác:} các chức năng chính đặt ở vị trí dễ quan sát, số bước thao tác được rút gọn.
    \item \textbf{Thống nhất:} màu sắc, kiểu nút, bảng dữ liệu và biểu mẫu dùng đồng nhất trên toàn hệ thống.
    \item \textbf{Phân tách theo vai trò:} mỗi nhóm người dùng chỉ thấy menu và chức năng phù hợp \texttt{role}.
    \item \textbf{Trực quan trạng thái:} trạng thái bàn, phiên phục vụ, món ăn và thanh toán hiển thị bằng nhãn/màu rõ ràng.
    \item \textbf{Phản hồi tức thì:} thông báo thành công/lỗi sau thao tác; màn bếp/phục vụ cập nhật realtime qua WebSocket.
    \item \textbf{Hỗ trợ mở rộng:} layout sidebar + router view cho phép bổ sung module mới mà không thay đổi khung tổng thể.
\end{itemize}

\subsection{Cấu trúc giao diện theo vai trò}

Bảng~\ref{tab:ui-roles} tóm tắt các nhóm màn hình chính của hệ thống theo từng vai trò.

\begin{table}[H]
\centering
\renewcommand{\arraystretch}{1.3}
\small
\begin{tabularx}{\textwidth}{|>{\RaggedRight\arraybackslash}p{3.2cm}|>{\RaggedRight\arraybackslash}X|>{\RaggedRight\arraybackslash}p{4.5cm}|}
\hline
\textbf{Vai trò} & \textbf{Chức năng giao diện} & \textbf{Màn hình tiêu biểu} \\
\hline
Khách hàng & Xem menu, đặt bàn, theo dõi bàn, đánh giá & \texttt{MenuView}, \texttt{Booking}, \texttt{MyTable}, \texttt{UserDashboard} \\ \hline
Nhân viên phục vụ & Quản lý bàn, tạo order, thanh toán & \texttt{AdminTables}, \texttt{Ordermenu} \\ \hline
Nhân viên bếp & Nhận và cập nhật trạng thái món & \texttt{AdminKitchen} \\ \hline
Quản lý chi nhánh & Dashboard, menu, báo cáo, đánh giá & \texttt{AdminDashboard}, \texttt{AdminReports}, \texttt{AdminReviews} \\ \hline
Admin hệ thống & Quản lý chi nhánh, tài khoản, nhân sự & \texttt{BranchManagement}, \texttt{CustomerAccounts}, \texttt{EmployeeManagement} \\ \hline
\end{tabularx}
\caption{Cấu trúc giao diện theo vai trò}
\label{tab:ui-roles}
\end{table}

\subsection{Các màn hình giao diện tiêu biểu}

Trong phạm vi báo cáo, nhóm trình bày các màn hình tiêu biểu đại diện cho các nghiệp vụ chính như sau.

\subsubsection{Giao diện đăng nhập}

Màn hình đăng nhập được trình bày như sau.

\begin{figure}[H]
    \centering
    \includegraphics[width=0.9\textwidth]{Chuong/UIDATN/login.png}
    \caption{Giao diện đăng nhập hệ thống}
    \label{fig:ui-login}
\end{figure}

Hình~\ref{fig:ui-login} hỗ trợ người dùng đăng nhập bằng tên đăng nhập và mật khẩu; sau khi xác thực thành công,
hệ thống cấp JWT và điều hướng giao diện theo vai trò tương ứng.

\subsubsection{Giao diện trang menu}

Màn hình menu dành cho khách hàng (UC04) được trình bày như sau.

\begin{figure}[H]
    \centering
    \includegraphics[width=\textwidth]{Chuong/menuUI.png}
    \caption{Giao diện trang menu dành cho khách hàng}
    \label{fig:ui-menu}
\end{figure}

Hình~\ref{fig:ui-menu} hiển thị danh sách món theo chi nhánh: tên món, mô tả, giá, ảnh và trạng thái còn bán;
khách có thể lọc theo danh mục trước khi đặt bàn hoặc gọi món tại bàn.

\subsubsection{Giao diện đặt bàn}

Màn hình đặt bàn trực tuyến (UC05) được trình bày như sau.

\begin{figure}[H]
    \centering
    \includegraphics[width=\textwidth]{Chuong/TableBookingUI.png}
    \caption{Giao diện đặt bàn trực tuyến}
    \label{fig:ui-booking}
\end{figure}

Hình~\ref{fig:ui-booking} cho phép chọn chi nhánh, ngày/giờ, số khách và ghi chú;
hệ thống kiểm tra bàn trống theo buffer 2 giờ, hiển thị gợi ý ``còn chỗ'' hoặc ``hết chỗ'' và áp dụng CAPTCHA chống spam.

\subsubsection{Giao diện quản lý bàn}

Màn hình quản lý bàn dành cho nhân viên phục vụ (UC06, UC10) được trình bày như sau.

\begin{figure}[H]
    \centering
    \includegraphics[width=\textwidth]{Chuong/waiterTable.png}
    \caption{Giao diện quản lý bàn và xác nhận đặt bàn}
    \label{fig:ui-waiter}
\end{figure}

Hình~\ref{fig:ui-waiter} hỗ trợ xem sơ đồ/trạng thái bàn, tiếp nhận khách walk-in, xác nhận phiên đặt trước,
xử lý no-show và chuyển bàn sang trạng thái phục vụ.

\subsubsection{Giao diện bếp}

Màn hình xử lý món ăn của bộ phận bếp (UC08) được trình bày như sau.

\begin{figure}[H]
    \centering
    \includegraphics[width=0.92\textwidth]{Chuong/UIDATN/kitchenUI.png}
    \caption{Giao diện xử lý món ăn của bộ phận bếp}
    \label{fig:ui-kitchen}
\end{figure}

Hình~\ref{fig:ui-kitchen} tối giản, ưu tiên danh sách món cần chế biến theo thứ tự thời gian;
nhân viên bếp cập nhật trạng thái \texttt{processing} $\rightarrow$ \texttt{done}, phục vụ nhận cập nhật realtime qua WebSocket.

\subsubsection{Giao diện quản lý chi nhánh}

Màn hình quản lý chi nhánh (UC14, UC15) được trình bày như sau.

\begin{figure}[H]
    \centering
    \includegraphics[width=0.92\textwidth]{Chuong/adminUI.png}
    \caption{Giao diện quản lý chi nhánh}
    \label{fig:ui-admin}
\end{figure}

Hình~\ref{fig:ui-admin} tập trung giám sát vận hành: thống kê nhanh doanh thu, số phiên phục vụ,
menu sidebar tới quản lý thực đơn, báo cáo và đánh giá khách hàng.

\section{Hiện thực các chức năng chính}
\label{section:4.4}

Trong phạm vi đồ án, nhóm đã hiện thực các chức năng cốt lõi của hệ thống quản lý chuỗi nhà hàng nhằm đáp ứng các nghiệp vụ chính trong quá trình vận hành thực tế.
Các màn hình tiêu biểu đã trình bày ở mục~\ref{section:4.3}; phần này mô tả cách hiện thực logic nghiệp vụ tương ứng.

\subsection{Chức năng đăng ký và đăng nhập}

Hệ thống hỗ trợ người dùng đăng ký và đăng nhập bằng tài khoản cá nhân (tên đăng nhập, mật khẩu, họ tên, số điện thoại).
Thông tin được lưu trên PostgreSQL; mật khẩu được mã hóa bằng bcrypt.
Form đăng ký tích hợp CAPTCHA và giới hạn tần suất gọi API nhằm hạn chế đăng ký tự động.

Sau khi đăng nhập thành công, hệ thống cấp JWT, phân quyền và điều hướng giao diện theo vai trò: khách hàng, nhân viên phục vụ, nhân viên bếp, quản lý chi nhánh hoặc admin hệ thống (Hình~\ref{fig:ui-login}).

\subsection{Chức năng xem menu và đặt bàn}

Khách hàng có thể xem danh sách món ăn theo từng chi nhánh trên Hình~\ref{fig:ui-menu}, bao gồm hình ảnh, mô tả, giá bán và trạng thái còn phục vụ.

Hệ thống hỗ trợ đặt bàn trực tuyến qua Hình~\ref{fig:ui-booking}: khách chọn chi nhánh, thời gian đến và số lượng khách;
hệ thống kiểm tra bàn trống trong khung giờ (tránh trùng lịch) và \textbf{tự động gán} bàn phù hợp theo sức chứa.
Khách có thể đặt món trước (pre-order) gắn với lịch đặt bàn.
Sau khi xác nhận, dữ liệu đồng bộ sang giao diện tiếp nhận và quản lý bàn của nhân viên phục vụ (Hình~\ref{fig:ui-waiter}).

\subsection{Chức năng tạo order và xử lý món ăn}

Nhân viên phục vụ tạo order món ăn cho bàn đang phục vụ (hoặc tiếp nhận khách walk-in/check-in đặt bàn trước) trên Hình~\ref{fig:ui-waiter}.
Sau khi order được tạo, danh sách món được gửi tới giao diện bếp (Hình~\ref{fig:ui-kitchen}).

Nhân viên bếp cập nhật trạng thái món ăn:
\begin{itemize}
    \item Chờ chế biến
    \item Đang chế biến
    \item Hoàn thành
\end{itemize}

Nhân viên phục vụ đánh dấu ``Đã phục vụ'' khi giao món tới bàn.
Các thay đổi trạng thái được đồng bộ qua \textbf{WebSocket} theo chi nhánh, kết hợp cơ chế làm mới định kỳ trên giao diện, giúp bếp và phục vụ theo dõi tiến độ gần thời gian thực.

\subsection{Chức năng thanh toán và xuất hóa đơn}

Hệ thống tổng hợp hóa đơn tạm tính theo các món đã gọi.
Khách hàng có thể \textbf{gửi yêu cầu thanh toán} trên màn hình ``Bàn của tôi'' hoặc xem hóa đơn qua mã QR tại bàn; mã VietQR hỗ trợ chuyển khoản, tuy nhiên việc \textbf{xác nhận thanh toán} do nhân viên phục vụ thực hiện trên hệ thống.

Nhân viên chọn phương thức: tiền mặt, chuyển khoản, thẻ hoặc ví điện tử, sau đó xác nhận và có thể xuất hóa đơn PDF.

Sau khi thanh toán thành công:
\begin{itemize}
    \item Đơn hàng và phiên phục vụ được cập nhật trạng thái đã thanh toán
    \item Bàn chuyển sang trạng thái ``Chờ dọn''
    \item Thông tin thanh toán được lưu phục vụ báo cáo doanh thu
\end{itemize}

\textbf{Lưu ý:} Backend có sẵn API tích hợp MoMo (môi trường thử nghiệm), nhưng luồng thanh toán end-to-end phía khách hàng chưa được đưa vào giao diện trong phạm vi đồ án.

\subsection{Chức năng quản lý và báo cáo}

Quản lý chi nhánh có thể trên Hình~\ref{fig:ui-admin}:
\begin{itemize}
    \item Xem và chỉnh sửa thông tin chi nhánh được phân công
    \item Quản lý thực đơn món ăn
    \item Xem báo cáo doanh thu, món bán chạy và xuất file Excel/PDF
    \item Xem đánh giá của khách hàng
\end{itemize}

Nhân viên phục vụ quản lý bàn, tiếp nhận khách, order và thanh toán trên cùng giao diện vận hành.
Admin hệ thống quản lý đa chi nhánh, nhân viên, tài khoản khách hàng và phân quyền.

\section{Kết quả kiểm thử}
\label{section:4.5}

Sau quá trình hiện thực, nhóm tiến hành \textbf{kiểm thử thủ công} (checklist theo từng use case và luồng nghiệp vụ) nhằm đảm bảo hệ thống hoạt động ổn định trong phạm vi đồ án.
Dự án chưa xây dựng bộ kiểm thử tự động (unit test/integration test).

\subsection{Kiểm thử chức năng}

Các chức năng chính đã được kiểm thử bao gồm:
\begin{itemize}
    \item Đăng ký và đăng nhập tài khoản (kèm CAPTCHA)
    \item Xem menu món ăn và danh sách chi nhánh
    \item Đặt bàn trực tuyến và đặt món trước
    \item Tiếp nhận khách, tạo order và cập nhật trạng thái món ăn
    \item Thanh toán, xuất hóa đơn và báo cáo doanh thu
    \item Phân quyền theo vai trò
\end{itemize}

Kết quả cho thấy các chức năng hoạt động đúng với luồng nghiệp vụ đã phân tích ở chương trước.

\subsection{Kiểm thử giao diện}

Giao diện được kiểm thử trên nhiều kích thước màn hình và trình duyệt:
\begin{itemize}
    \item Google Chrome
    \item Microsoft Edge
    \item Mozilla Firefox
\end{itemize}

Kết quả cho thấy giao diện hiển thị ổn định, thao tác dễ sử dụng và đảm bảo tính trực quan cho người dùng.

\subsection{Kiểm thử hiệu năng}

Hệ thống được thử nghiệm với các thao tác lặp lại và nhiều thao tác nối tiếp như đặt bàn, gọi món, cập nhật trạng thái món ăn và thanh toán trong điều kiện dữ liệu mức đồ án (một số chi nhánh, vài chục tài khoản mẫu).

Thời gian phản hồi trung bình của các thao tác phổ biến thường dưới 1 giây trên môi trường phát triển local; chưa đánh giá tải với hàng trăm người dùng đồng thời.

\subsection{Đánh giá kết quả kiểm thử}

Thông qua kiểm thử thủ công, hệ thống đáp ứng các yêu cầu chức năng và phi chức năng cơ bản của đề tài.
Dữ liệu đồng bộ ổn định giữa phục vụ, bếp và quản lý, góp phần giảm sai sót trong vận hành nhà hàng.

\section{Đánh giá hệ thống}
\label{section:4.6}

\subsection{Kết quả đạt được}

Hệ thống đã xây dựng thành công các chức năng cốt lõi của quản lý chuỗi nhà hàng nhiều chi nhánh.

Các chức năng chính như:
\begin{itemize}
    \item Quản lý tài khoản và phân quyền
    \item Xem menu, đặt bàn và đặt món trước
    \item Tạo order, xử lý món ăn và đồng bộ bếp--phục vụ
    \item Thanh toán do nhân viên xác nhận, xuất hóa đơn
    \item Báo cáo doanh thu và quản lý đa chi nhánh
\end{itemize}

đều đã được hiện thực và hoạt động tương đối ổn định trong phạm vi nguyên mẫu.

Hệ thống áp dụng kiến trúc phân tầng kết hợp mô hình MVC, thuận tiện mở rộng và bảo trì.

\subsection{Ưu điểm của hệ thống}

Hệ thống có một số ưu điểm nổi bật:
\begin{itemize}
    \item Giao diện trực quan, phân tách rõ theo vai trò người dùng.
    \item Phân quyền theo vai trò và phạm vi chi nhánh (đối với quản lý).
    \item Đồng bộ trạng thái món ăn qua WebSocket, hỗ trợ vận hành bếp--phục vụ.
    \item Hỗ trợ QR bàn, VietQR và xuất báo cáo Excel/PDF.
    \item Kiến trúc tách biệt frontend--backend, dễ mở rộng.
\end{itemize}

\subsection{Hạn chế của hệ thống}

Bên cạnh các kết quả đạt được, hệ thống vẫn còn hạn chế:
\begin{itemize}
    \item Chưa tích hợp cổng thanh toán điện tử end-to-end phía khách hàng (MoMo mới ở mức API thử nghiệm, chưa gắn giao diện).
    \item VietQR chỉ hỗ trợ hiển thị mã chuyển khoản; chưa tự động đối soát thanh toán.
    \item Chưa tối ưu cho số lượng người dùng đồng thời rất lớn; chưa kiểm thử tải chính thức.
    \item WebSocket mới phục vụ đồng bộ nội bộ chi nhánh, chưa mở rộng toàn hệ thống phức tạp.
    \item Chưa có ứng dụng di động riêng.
\end{itemize}

\subsection{Định hướng cải tiến}

Trong tương lai, hệ thống có thể phát triển thêm:
\begin{itemize}
    \item Hoàn thiện thanh toán MoMo/VNPay trên giao diện khách hàng.
    \item Tự động đối soát chuyển khoản và tối ưu kênh WebSocket.
    \item Bổ sung module quản lý kho và khuyến mãi gắn hóa đơn.
    \item Bổ sung kiểm thử tự động và đánh giá hiệu năng dưới tải cao.
    \item Ứng dụng AI hỗ trợ dự báo doanh thu và lượng khách.
\end{itemize}

\end{document}
